\section{Conclusion}
\label{sec:conclusion}

HMRS manages documents at three levels and answers questions along a path that
mirrors them. At the collection level, documents are parsed, classified into
the library whose scope matches their body, and catalogued with derived
descriptions and tags; at the document level, move, rename, delete and forced
re-index operations update the collection, while reconciliation can detect
selected disagreements among three state sources; at the segment level, an
optional gate scores candidate passages and returns survivors with provenance.
The collection-selection step connects the managed catalogue to retrieval;
its outcome depends on the chosen access workflow.

The campaigns confirm selected paths of that structure and also expose its
limits. Ingestion and classification produced 322 stored parts across five
libraries, but 21 of 100 submissions were not retrievable and the surviving
records cannot attribute the failure to a stage. A separate 33-part storage
sample passed, while an independent disposable-library test exposed a
stale-index rename defect corrected by forced re-indexing. In 16 lifecycle
scenarios, create, move, delete and forced re-index passed, whereas task-status
reporting and the update check failed. The route-card audit found
collection-selection correctness of 72/72 with full route cards versus 37/72
with name-only cards, but an unlogged fallback
rules out a causal interpretation. Evidence-pack string containment was
37/72 (A0) versus 33/72 (A3) ($p=0.4545$). A mechanical full-read script found all 72
strings after reading 879 parts, but generated no answers. The answer grids had
gold identifiers in the answer or recorded trace for all 60 cells, but lexical and citation markers do not
establish semantic support. On identical pools, verifier re-ranking reduced
Hit@5 from 0.94 to 0.44; fusion order restored 0.94, and the configured
threshold rejected none of 191 scores. On the same 72-question corpus, the
tiered path with a chunk-window evidence union contained the gold string in
52/72, above a scoped flat union with a potentially gold-dependent reference
lookup (39/72; 13--0, unadjusted $p=0.00024$) and above a precise arm without navigation (46/72;
7--1, $p=0.0703$). The answer grids add a substrate contrast: answers over the
document-level organisation carried document-style citation markers in nine
of ten cells against two of ten over a flat 800-character chunk replica of
the same content, a difference bound to the content substrate rather than a
retrieval difference. A recorded demonstration on one newly ingested paper showed two
agent harnesses passing the same four-question case review on content,
provenance, quotation and absence of fabrication (supplementary Table~S14),
including a matched refusal to report a cost figure the source does not state.

These results document an auditable management structure, bounded per-tier
measurements and a deployment warning against verifier-as-ranker. The
content-layer comparisons change multiple components, and the flat arm uses a
potentially gold-label-dependent lookup; they do not isolate a hierarchy effect. They
do not establish system-level superiority, general retrieval effectiveness,
semantic answer quality or reliability. Broader evaluation needs independently
authored questions, repeated agent trials with run-to-run variance estimates,
claim-level support review, larger and repeated versions of the
ladder-versus-flat-and-precise contrast with multiplicity control, and
crash-recovery tests. Within these bounds, HMRS shows that treating the
collection as a managed, auditable object is a viable, inspectable basis for
evidence-traced retrieval.
